% Figure: the linear-extrapolation trap on an exponential process.
% Early exponential data looks nearly linear; a linear fit through it
% and an exponential fit through it agree in the fitted window and
% diverge by orders of magnitude outside it.
\begin{figure}[htbp]
\centering
\begin{tikzpicture}[scale=1.0,
  ax/.style={->, draw=engblack, thick},
  tick/.style={draw=engblack!60, thin},
  expo/.style={draw=engred, thick},
  linf/.style={draw=engblue, thick, dashed},
  fitwin/.style={draw=engamber, thin},
  pt/.style={circle, fill=engblack, inner sep=1.2pt},
  ann/.style={font=\footnotesize, align=left},
]
% --- axes ---
\draw[ax] (0, 0) -- (6.5, 0) node[below right, ann] {$t$};
\draw[ax] (0, 0) -- (0, 5.3) node[left, ann] {$y$};
\foreach \x in {1, 2, 3, 4, 5, 6} {
  \draw[tick] (\x, -0.05) -- (\x, 0.05);
  \node[ann, below, yshift=-2pt] at (\x, 0) {$\x$};
}
% fitted window 0..2 shaded by a bracket
\draw[fitwin] (0, -0.25) -- (2, -0.25);
\draw[fitwin] (0, -0.18) -- (0, -0.32);
\draw[fitwin] (2, -0.18) -- (2, -0.32);
\node[ann, anchor=north] at (1, -0.35) {fitted window};
% exponential y = 0.3 e^{0.8 t}, scaled to fit
\draw[expo, domain=0:3.05, samples=50, smooth]
  plot (\x, {0.3*exp(0.8*\x)});
\node[ann, anchor=south east, draw=engred, fill=white, inner sep=2pt]
  at (3.0, 4.9) {true: $y = 0.3 e^{0.8t}$};
% data points in window
\foreach \x in {0, 0.5, 1.0, 1.5, 2.0} {
  \pgfmathsetmacro{\yy}{0.3*exp(0.8*\x)}
  \node[pt] at (\x, \yy) {};
}
% linear fit through the window: slope from (0,0.3) to (2, 1.49)
% slope = (1.49-0.3)/2 = 0.595; line y = 0.3 + 0.595 t
\draw[linf, domain=0:6.3, samples=2]
  plot (\x, {0.3 + 0.595*\x});
\node[ann, anchor=west, draw=engblue, fill=white, inner sep=2pt]
  at (4.2, 2.7) {linear extrapolation};
\end{tikzpicture}
\caption[The linear-extrapolation trap on an exponential]{An
exponential process $y = 0.3 e^{0.8 t}$ (red) sampled over a fitted
window $t \in [0, 2]$ (black points). Across the window the curve is
nearly straight, and a linear fit (blue dashed) tracks it closely; an
analyst who saw only the windowed data would not distinguish the two.
Outside the window the curves diverge without bound: at $t = 3$ the
exponential is already climbing past the top of the frame while the
linear extrapolation predicts a modest value. The visual closeness
inside the fitted window is exactly what makes the linear
extrapolation seductive and wrong, and the remedy is to plot the
process on semilog axes where the exponential is the straight line
and the extrapolation is honest.}
\label{fig:vol02:ch02:extrapolation-trap}
\end{figure}
